% SD4H Methods — TARGET: ~1.0 page in ICML two-column format

\section{Methods}
\label{sec:method}

\subsection{Problem Setting}
\label{sec:problem}

Let $\mathcal{X}_S \sim P_S$ denote source-domain clinical time series (\eicu{} or \hirid{}) and $\mathcal{X}_T \sim P_T$ denote the target domain (\mimic{}).
A predictor $f_T : \mathbb{R}^{T \times F} \to \mathbb{R}^{T \times C}$, trained on $\mathcal{X}_T$, is \emph{entirely frozen}: all parameters satisfy $\nabla_\phi f_T = 0$.
We learn an adapter $g_\theta : \mathbb{R}^{T \times F} \to \mathbb{R}^{T \times F}$ that maps source inputs into the target input space so that the composite system $f_T \circ g_\theta$ performs well on source labels $y_S$:
\begin{multline}
\label{eq:objective}
\mathcal{L}(\theta) = \mathcal{L}_{\mathrm{task}}\!\bigl(f_T(g_\theta(x_S)),\, y_S\bigr) \\
+ \lambda_{\mathrm{fid}}\,\mathcal{L}_{\mathrm{fid}}\!\bigl(g_\theta(x_S),\, x_S\bigr)
+ \lambda_{\mathrm{range}}\,\mathcal{L}_{\mathrm{range}}\!\bigl(g_\theta(x_S)\bigr),
\end{multline}
where $\mathcal{L}_{\mathrm{task}}$ is cross-entropy (classification) or MSE (regression) evaluated through the frozen predictor, $\mathcal{L}_{\mathrm{fid}}$ is an MSE fidelity penalty regularizing toward identity, and $\mathcal{L}_{\mathrm{range}}$ penalizes outputs outside the physiological range observed in $\mathcal{X}_T$.
\Cref{eq:objective} governs Phase~2 (adaptation training); Phase~1 uses a separate reconstruction objective on target data only (\S\ref{sec:training}).
The frozen constraint reflects deployment reality: under the FDA SaMD framework, validated predictors cannot be modified post-certification~\citep{fda2021samd}.

\subsection{Retrieval-Guided Architecture}
\label{sec:architecture}

% Architecture diagram — Retrieval-Guided Adapter
% Compact top-to-bottom layout for single-column ICML format
\begin{figure}[t]
\centering
\resizebox{0.7\columnwidth}{!}{%
\begin{tikzpicture}[
  % Global styles
  >=Stealth,
  font=\small,
  % Block styles
  trainable/.style={
    draw=blue!70!black, fill=blue!8, rounded corners=3pt,
    minimum height=0.7cm, minimum width=2.6cm,
    font=\small\bfseries, text=blue!70!black,
    line width=0.8pt
  },
  frozen/.style={
    draw=red!60!black, fill=red!6, rounded corners=3pt,
    minimum height=0.7cm, minimum width=2.6cm,
    font=\small\bfseries, text=red!60!black,
    line width=1.2pt, double, double distance=0.8pt
  },
  membank/.style={
    draw=green!50!black, fill=green!8,
    cylinder, shape border rotate=90,
    cylinder uses custom fill, cylinder end fill=green!15,
    cylinder body fill=green!8,
    minimum height=1.15cm, minimum width=2.0cm,
    font=\small\bfseries, text=green!50!black,
    line width=0.8pt, inner sep=3pt, align=center
  },
  databox/.style={
    draw=black!50, fill=gray!5, rounded corners=2pt,
    minimum height=0.55cm, minimum width=1.8cm,
    font=\small, text=black!70,
    line width=0.6pt
  },
  smallbox/.style={
    draw=blue!70!black, fill=blue!8, rounded corners=2pt,
    minimum height=0.5cm, minimum width=1.6cm,
    font=\scriptsize\bfseries, text=blue!70!black,
    line width=0.6pt
  },
  arr/.style={->, line width=0.7pt, black!70},
  darr/.style={->, line width=0.7pt, green!50!black, dashed}
]

% =====================================================
% Main vertical flow (left-center column)
% =====================================================

% Source input
\node[databox] (src) at (0, 0) {Source input $x_S$};

% Normalize
\node[smallbox] (norm) at (0, -0.85) {Normalize};

% Encoder
\node[trainable] (enc) at (0, -1.65) {Shared Encoder $E_\phi$};

% Cross-Attention
\node[trainable] (cross) at (0, -2.6) {Cross-Attention $\times\!L$};

% Decoder
\node[trainable] (dec) at (0, -3.5) {Decoder $D_\psi$};

% Adapted output
\node[databox] (out) at (0, -4.45) {Adapted $\tilde{x}$};

% Frozen predictor
\node[frozen] (pred) at (0, -5.35) {Frozen Model $f_T$};

% Prediction output
\node[font=\small\bfseries, text=black!70] (ypred) at (0, -6.0) {Prediction $\hat{y}$};

% =====================================================
% Memory bank (right side)
% =====================================================

% Position memory bank between encoder and cross-attention level
\node[membank] (mem) at (3.2, -1.7) {%
  Target\\[-1pt]Memory\\[-1pt]Bank $\mathcal{M}$};

% Target data feeding memory bank
\node[databox] (tgt) at (3.2, -0.4) {Target data $\mathcal{X}_T$};

% =====================================================
% Arrows -- main vertical flow
% =====================================================
\draw[arr] (src) -- (norm);
\draw[arr] (norm) -- (enc);
\draw[arr] (enc) -- (cross);
\draw[arr] (cross) -- (dec);
\draw[arr] (dec) -- (out);
\draw[arr] (out) -- (pred);
\draw[arr] (pred) -- (ypred);

% =====================================================
% Arrows -- memory bank interactions
% =====================================================

% Target data -> memory bank (dashed green)
\draw[darr] (tgt) -- (mem);

% Encoder -> memory bank: horizontal arrow, label above
\draw[arr, blue!60!black]
  (enc.east) -- (mem.west)
  node[midway, above, font=\scriptsize, text=black!60] {query};

% Memory bank -> cross-attention: down from bank center, then left
% Place waypoint directly below memory bank center at cross-attention y-level
\coordinate (knnTurn) at (3.2, -2.6);
\draw[arr, green!50!black]
  (mem.south) -- (knnTurn) -- (cross.east);
% k-NN label on horizontal segment, midway between turn and cross-attn
\node[above, font=\scriptsize, text=black!60]
  at ($(knnTurn)!0.5!(cross.east)$) {$k$-NN};

% =====================================================
% Adapter bounding box (explicit rectangle)
% =====================================================
\begin{scope}[on background layer]
  \draw[draw=blue!35, dashed, rounded corners=6pt, line width=0.8pt, fill=blue!2]
    ([xshift=-0.95cm, yshift=0.1cm] norm.north west)
    rectangle
    ([xshift=0.4cm, yshift=-0.28cm] dec.south east);
\end{scope}
% Label in the expanded whitespace below the decoder, inside the box
\node[font=\tiny\bfseries, text=blue!50!black, anchor=north east, inner sep=0.5pt]
  at ([xshift=0.3cm, yshift=-0.02cm] dec.south east)
  {Adapter $g_\theta$};

% =====================================================
% Frozen indicator -- snowflake symbol to the left
% =====================================================
\begin{scope}[shift={([xshift=-0.55cm] pred.west)}, scale=0.5]
  % Six-armed snowflake
  \foreach \a in {0, 60, 120} {
    \draw[red!60!black, line width=0.7pt] (\a:0.35) -- (\a+180:0.35);
    % Small branches on each arm
    \draw[red!60!black, line width=0.5pt] (\a:0.2) -- ++(\a+45:0.1);
    \draw[red!60!black, line width=0.5pt] (\a:0.2) -- ++(\a-45:0.1);
    \draw[red!60!black, line width=0.5pt] (\a+180:0.2) -- ++(\a+180+45:0.1);
    \draw[red!60!black, line width=0.5pt] (\a+180:0.2) -- ++(\a+180-45:0.1);
  }
  \fill[red!60!black] (0,0) circle (0.04);
\end{scope}

\end{tikzpicture}%
}% end resizebox
\caption{The \retr{} architecture. Source input is encoded, matched to target exemplars by $k$-NN, fused through cross-attention, and decoded for the frozen downstream model.}
\label{fig:arch}
\end{figure}


The \retr{} architecture (\Cref{fig:arch}) comprises four components: a shared encoder, a target-domain memory bank, cross-attention fusion blocks, and a decoder.

\paragraph{Shared encoder and memory bank.}
A shared encoder $E_\phi$ maps both source and target windows to a latent space of dimension $d_{\text{latent}}$: $z_s = E_\phi(x_s) \in \mathbb{R}^{T \times d_{\text{latent}}}$.
Each encoder and decoder layer is an \emph{axial attention block} that applies self-attention along both temporal and feature dimensions, so the model handles per-feature dynamics and cross-feature interactions without concatenating the full $T \times F$ sequence.
All target-domain training data is encoded and segmented into temporal windows of size $W$ timesteps, each mean-pooled into a compact descriptor (Appendix~\ref{app:implementation}).
The resulting memory bank $\mathcal{M} = \{E_\phi(x_T^j)\}_{j=1}^{|\mathcal{X}_T|}$ resides on GPU for fast retrieval and is rebuilt every $R$ epochs as the encoder's representations evolve.

\paragraph{Per-timestep $k$-NN retrieval.}
For each source timestep $t$, a query is formed by mean-pooling source latents over a backward-looking window $[\max(0, t{-}W{+}1),\, t]$, preserving causal consistency.
Learned importance weights $\sigma(\mathbf{w}) \in [0,1]^{d_{\text{latent}}}$ modulate a weighted Euclidean distance over latent dimensions, and the $k$ nearest target windows are retrieved, with $k$ selected by grid search (Appendix~\ref{app:implementation}).
This per-timestep design matches local physiological states rather than full trajectories and outperforms wider retrieval windows in all tasks.

\paragraph{Cross-attention fusion.}
Retrieved target context is integrated through $L$ stacked cross-attention blocks.
Let $\mathcal{N}_k(z_s^t)$ denote the $k$ nearest target latents for timestep $t$; each block computes:
%
\begin{equation}
\label{eq:cross}
z_{\mathrm{cross}}^t = \mathrm{CrossAttn}\!\bigl(Q{=}z_s^t,\; K{=}V{=}\mathcal{N}_k(z_s^t)\bigr) + z_s^t,
\end{equation}
%
followed by causal self-attention and a feed-forward network, both with residual connections and pre-LayerNorm.
Cross-attention allows selective, feature-level borrowing from target exemplars, so the model attends to different features from different neighbors rather than uniformly interpolating as in $k$NN-MT~\citep{khandelwal2021knnmt}.
The decoder $D_\psi$ maps the fused representations back to input space:
%
\begin{equation}
\label{eq:decode}
\tilde{x} = D_\psi(z_{\mathrm{cross}}) \in \mathbb{R}^{T \times F}.
\end{equation}
%
The cross-attention depth $L$ is selected per task (Appendix~\ref{app:implementation}).

\subsection{Training}
\label{sec:training}

Training follows a two-phase protocol.
\textbf{Phase~1} (autoencoder pretrain): $E_\phi$ and $D_\psi$ are trained on target data $\mathcal{X}_T$ only with a reconstruction objective; ablations (Appendix~\ref{app:ablation}) show this phase is critical.
\textbf{Phase~2} (adaptation): end-to-end training on source data through the frozen predictor with the objective of \Cref{eq:objective}, where $\mathcal{L}_{\mathrm{task}}$ supplies the primary gradient through $f_T$ and $\mathcal{L}_{\mathrm{fid}}$, $\mathcal{L}_{\mathrm{range}}$ regularize the output; the memory bank is rebuilt every $R$ epochs as representations evolve.

Before adaptation, source features undergo cross-domain normalization: an affine transform matching per-feature mean and variance to target statistics, computed once from training data.
Attention is causal for per-timestep tasks (AKI, Sepsis, LoS) and bidirectional for per-stay tasks (Mortality, KF), preventing information leakage across time. AdaTime datasets (HAR, HHAR, WISDM, SSC, MFD) present fixed-length windows with a single label per window and use bidirectional attention throughout.
